\documentclass[12pt]{article}
\def\activityid{F09-M-v3}\usepackage[letterpaper,margin=.65in,top=.60in,bottom=.65in]{geometry}
\usepackage[T1]{fontenc}
\usepackage[default]{sourcesanspro}
\usepackage{microtype,tikz,array,tabularx,fancyhdr,amsmath,amssymb}
\usepackage[hidelinks]{hyperref}
\usetikzlibrary{calc,arrows.meta}
\setlength{\parindent}{0pt}\setlength{\parskip}{7pt}
\setlength{\headheight}{0pt}\setlength{\footskip}{26pt}\setlength{\emergencystretch}{2em}
\pagestyle{fancy}\fancyhf{}\renewcommand{\headrulewidth}{0pt}\renewcommand{\footrulewidth}{.4pt}
\fancyfoot[L]{\scriptsize Bellingham Math Circle / Week 9 / \activityid}
\fancyfoot[R]{\scriptsize \thepage}
\definecolor{ink}{gray}{.12}\definecolor{soft}{gray}{.95}\color{ink}
\newcommand{\PageTitle}[2]{{\fontsize{10}{12}\selectfont\bfseries WEEK 9\quad BOUNCE LAB\hfill #1\par}\vspace{3pt}{\fontsize{26}{29}\selectfont\bfseries #2\par}\vspace{2pt}}
\newcommand{\Subhead}[1]{\par\vspace{3pt}{\large\bfseries #1}\par}
\newcommand{\RuleBox}[1]{\par\begingroup\setlength{\fboxsep}{8pt}\colorbox{soft}{\parbox{\dimexpr\linewidth-16pt\relax}{#1}}\endgroup\par}
\newcommand{\WriteBox}[2]{\par\textbf{#1}\par\vspace{2pt}\begin{tikzpicture}\draw[black!40,line width=.5pt] (0,0) rectangle (\dimexpr\linewidth-.5pt\relax,#2);\end{tikzpicture}\par}
\newcommand{\RookBoard}[3]{\begin{tikzpicture}[x=#3,y=#3]\draw[step=1,black!45] (0,0) grid (#1,#2);\draw[thick] (0,0) rectangle (#1,#2);\node[font=\Large] at ({#1-.5},.5){$\star$};\draw[-{Latex},thick] (0,-.35)--(#1,-.35);\draw[-{Latex},thick] ({#1+.35},#2)--({#1+.35},0);\end{tikzpicture}}
\newcommand{\Table}[3]{\begin{tikzpicture}[x=#3,y=#3]\draw[step=1,black!25] (0,0) grid (#1,#2);\draw[thick] (0,0) rectangle (#1,#2);\fill (0,0)circle(3pt);\draw[-{Latex},very thick] (0,0)--(.7,.7);\node[below] at ({#1/2},-.1){width #1};\node[rotate=90,above] at (-.2,{#2/2}){height #2};\end{tikzpicture}}
\newcommand{\GraphNodes}{\coordinate(A)at(0,0);\coordinate(B)at(3,0);\coordinate(C)at(1.5,2.6);\coordinate(D)at(1.5,.95);}
\newcommand{\Kfour}[1]{\begin{tikzpicture}[scale=#1]\GraphNodes\draw[very thick](A)--(B)--(C)--(A)--(D)--(B) (D)--(C);\foreach \n in {A,B,C,D}{\node[circle,draw,fill=white,minimum size=22pt,font=\bfseries]at(\n){\n};}\end{tikzpicture}}
\newcommand{\House}[1]{\begin{tikzpicture}[scale=#1]\coordinate(A)at(0,0);\coordinate(B)at(3,0);\coordinate(C)at(3,2);\coordinate(D)at(0,2);\coordinate(E)at(1.5,3.3);\draw[very thick](A)--(B)--(C)--(D)--(A) (D)--(E)--(C);\foreach \n in {A,B,C,D,E}{\node[circle,draw,fill=white,minimum size=22pt,font=\bfseries]at(\n){\n};}\end{tikzpicture}}




\newcommand{\StudentPage}[1]{{\fontsize{10}{12}\selectfont\bfseries Week 9 / Billiards lab / #1\par}\vspace{10pt}}
\newcommand{\Problem}[2]{\par\noindent\textbf{Problem #1:} #2\par}
\newcommand{\Space}[1]{\par\begin{tikzpicture}\draw[black!35] (0,0) rectangle (\dimexpr\linewidth-.5pt\relax,#1);\end{tikzpicture}\par}
\newcommand{\Piles}[1]{\begin{center}\begin{tikzpicture}[x=1in,y=1in]\foreach \x in {0,...,#1}{\draw[black!50,rounded corners] ({\x*2.05},0) rectangle ({\x*2.05+1.8},1.3);}\end{tikzpicture}\end{center}}

\newcommand{\Rooms}[5]{\begin{tikzpicture}[x=#5,y=#5]\draw[step=1,black!20](0,0)grid(#3,#4);\foreach\x in{0,#1,...,#3}{\draw[thick](\x,0)--(\x,#4);}\foreach\y in{0,#2,...,#4}{\draw[thick](0,\y)--(#3,\y);}\fill(0,0)circle(3pt);\end{tikzpicture}}
\newcommand{\PlainGrid}[3]{\begin{tikzpicture}[x=#3,y=#3]\draw[step=1,black!30](0,0)grid(#1,#2);\fill(0,0)circle(3pt);\end{tikzpicture}}

\begin{document}

\StudentPage{Grades 2--3}
\Problem{1}{Start at each black dot and trace diagonally up-right, one grid step across for each step up. At a side wall reverse left/right; at a top or bottom wall reverse up/down. Stop at the first corner. Line crossings inside the table do not change the direction. Mark every bounce with a dot and circle the ending corner.}
\begin{center}\Table{2}{2}{.55in}\hspace{.5in}\Table{2}{4}{.55in}\end{center}
\begin{center}\Table{4}{2}{.55in}\end{center}
\Problem{2}{For each table, record the corner and number of bounces. Count wall bounces, not the starting or ending corner.}
\begin{center}\renewcommand{\arraystretch}{2}\begin{tabular}{c|p{2in}|p{2in}}Width $\times$ height&Ending corner&Bounces\\\hline$2\times2$&&\\$2\times4$&&\\$4\times2$&&\end{tabular}\end{center}
\newpage
\StudentPage{Grades 2--3}
\Problem{3}{Trace this $2\times3$ table. Mark bounces with dots. Record their order with small numbers beside the dots. Circle the ending corner.}
\begin{center}\Table{2}{3}{.49in}\end{center}
\Problem{4}{The thick lines below are walls of mirror copies of that table. Instead of bouncing, draw a straight diagonal from the black dot. Stop at the first meeting of a thick vertical and thick horizontal wall. Number the wall crossings before it. Match those numbers to your bounces in Problem 3.}
\begin{center}\Rooms{2}{3}{6}{6}{.43in}\end{center}
\Problem{5}{Count the whole room widths and heights from the black dot to your stopping point. Record: widths \rule{.6in}{.4pt}; heights \rule{.6in}{.4pt}; horizontal distance \rule{.6in}{.4pt}.}
\newpage
\StudentPage{Grades 2--3}
\Problem{6}{Draw a straight diagonal to the first meeting of a vertical and horizontal room wall in each picture. The left rooms have width 3, height 6. The right rooms have width 3, height 9. Mark each wall crossing. Count whole widths and heights to the stopping point.}
\begin{center}\Rooms{3}{6}{6}{6}{.31in}\hspace{.6in}\Rooms{3}{9}{9}{9}{.31in}\end{center}
\begin{center}\renewcommand{\arraystretch}{2}\begin{tabular}{c|c|c|p{1.7in}}Table&Widths&Heights&Ending corner\\\hline$3\times6$&&&\\$3\times9$&&&\end{tabular}\end{center}
\Problem{7}{Fold the trip back into its original room. One whole width carries you to the right; the next brings you back left. One whole height carries you to the top; the next brings you back down. Move your finger back and forth along these edges to find the ending side after each count. Use these records to finish the corner column above.}
\begin{center}\renewcommand{\arraystretch}{1.7}\begin{tabular}{l|*{6}{>{\centering\arraybackslash}p{.53in}}}Whole widths or heights&1&2&3&4&5&6\\\hline Left/right endpoint&right&&&&&\\Bottom/top endpoint&top&&&&&\end{tabular}\end{center}
\newpage
\StudentPage{Grades 2--3}
\Problem{8}{Use lists of wall distances to predict a corner. For width 4, list 4,8,12,\ldots. For height 6, list 6,12,18,\ldots. Stop at their first common number. Do the same for each row. Record whole room counts and use Problem 7 to predict the corner.}
\begin{center}\renewcommand{\arraystretch}{2}\begin{tabular}{c|p{1.4in}|c|c|p{1.15in}}Table&\shortstack{First common\\distance}&Widths&Heights&Corner\\\hline$4\times6$&&&&\\$2\times6$&&&&\\$3\times3$&&&&\end{tabular}\end{center}
\Problem{9}{On the left grid draw a rectangle that reaches the top-right corner after at least one bounce. On the right draw one that reaches the top-left. Use whole-number widths and heights no larger than 6. Trace each path to test your design.}
\begin{center}\PlainGrid{6}{6}{.40in}\hspace{.5in}\PlainGrid{6}{6}{.40in}\end{center}
\Problem{10}{Record each table's width, height, and ending corner. Swap designs with a partner and test their prediction.}
\Space{.6in}

\end{document}
